%% ma: L12-B14-0054
%% lop: 12
%% chuong: 5
%% bai: 14
%% dang: 12.14.1
%% loai: TN4
%% muc: NB
%% dap_an: "D"
%% nguon: "To_n_12_PT_M_t_Ph_ng.pdf (Đề kiểm tra đánh giá kiến thức, Chuyên đề Phương trình mặt phẳng, Nguyễn Anh Khôi - Lê Minh Khôi), Câu 8"
%% kien_thuc: "Bài 14 (phương trình tổng quát của mặt phẳng, vectơ pháp tuyến)"
%% trang_thai: chinh-thuc
%% ghi_chu: "Tài liệu gốc không kèm đáp án; đáp án do Claude giải và kiểm bằng sympy."
\begin{ex}
Trong không gian $Oxyz$, cho mặt phẳng $(P)\colon x-2y+3z+2020=0$. Vectơ nào dưới đây \textbf{không phải} là một vectơ pháp tuyến của mặt phẳng $(P)$?
\choice{$\vec n=(-2;4;-6)$}{$\vec n=(-1;2;-3)$}{$\vec n=(1;-2;3)$}{\True $\vec n=(-2;3;2020)$}
\loigiai{$(P)$ có vectơ pháp tuyến $(1;-2;3)$; các vectơ ở A, B, C cùng phương với vectơ này; vectơ $(-2;3;2020)$ không cùng phương. Chọn D.}
\end{ex}
